% ============================================================
%  SECCION 1 — GUION POR DIAPOSITIVA
% ============================================================
\section{Guion por Diapositiva}

\begin{tcolorbox}[colback=ujatblue!10, colframe=ujatblue,
  title={\bfseries Informacion general de la presentacion}, coltitle=white]
\begin{tabular}{@{}ll@{}}
\textbf{Duracion objetivo:} & $\leq$ 20 minutos (margen recomendado: 18 min) \\
\textbf{Total de diapositivas:} & 33 (incluyendo 3 de respaldo --- no contar) \\
\textbf{Diapositivas de contenido:} & 30 (S01--S30) \\
\textbf{Nota sobre animaciones:} & S16/S17, S18/S19, S20--S23 son pares de animacion. \\
& Presentar siempre el \textbf{ultimo} de cada grupo (tiene todo el contenido). \\
\end{tabular}
\end{tcolorbox}

% --- Resumen de tiempos ---
\subsection{Resumen de tiempos por bloque}

\begin{table}[h!]
\centering
\begin{tabular}{@{}llcr@{}}
\toprule
\textbf{Bloque} & \textbf{Slides} & \textbf{Tema} & \textbf{Tiempo} \\
\midrule
Bloque I  & S01--S09 & Introduccion, Problema, Marco Teorico & $\approx$ 5:30 min \\
Bloque II & S10--S15 & Marco Fisico, PINNs, Arquitectura, Metodologia & $\approx$ 4:30 min \\
Bloque III & S16--S24 & Resultados NB01 y NB02 & $\approx$ 5:30 min \\
Bloque IV & S25--S30 & Comparativa, Discusion, Cierre & $\approx$ 2:30 min \\
\midrule
\textbf{Total} & & & $\approx$ \textbf{18:00 min} \\
\bottomrule
\end{tabular}
\caption{Distribucion de tiempo por bloque. Margen de $\approx$2 min para imprevistos.}
\end{table}

\vspace{6pt}
\hrule
\vspace{12pt}

% ==============================================================
%  DIAPOSITIVA 1
% ==============================================================
\subsection{S01 --- Portada}

\begin{timebox}
Tiempo estimado: 0:25 min \quad|\quad Acumulado: 0:25 min
\end{timebox}
\vspace{4pt}

\begin{slidebox}{Guion}
``Buenos dias. Mi nombre es Roberto Hernandez Estrada, estudiante de Maestria en
Ciencias de la Computacion en la UJAT. El tema de este segundo avance de tesis es
\emph{la simulacion acelerada de speckle optico mediante Redes Neuronales Informadas
por Fisica con activacion sinusoidal}. Mi director es el Dr. Jose Adan Hernandez
Nolasco. Gracias por su tiempo.''
\end{slidebox}
\vspace{4pt}

\begin{tipbox}
Hablar de pie con contacto visual con cada sinodal. \textbf{No leer} la diapositiva.
Voz clara y volumen apropiado desde el inicio.
\end{tipbox}

\begin{warningbox}
Apoyo en diapositiva: titulo, institucion, nombre del alumno y directores.
\textbf{Solo senalar, no leer.}
\end{warningbox}

\vspace{8pt}

% ==============================================================
%  DIAPOSITIVA 2
% ==============================================================
\subsection{S02 --- Contenido de la Presentacion}

\begin{timebox}
Tiempo estimado: 0:25 min \quad|\quad Acumulado: 0:50 min
\end{timebox}
\vspace{4pt}

\begin{slidebox}{Guion}
``La presentacion esta organizada en tres bloques. El \textbf{Bloque I} cubre la
introduccion: planteamiento del problema, preguntas de investigacion, hipotesis
y marco teorico. El \textbf{Bloque II} presenta la metodologia y los resultados
experimentales de los dos primeros notebooks validados. El \textbf{Bloque III}
cierra con la comparativa contra el estado del arte, los hallazgos del avance y
los proximos pasos.''
\end{slidebox}
\vspace{4pt}

\begin{tipbox}
Senalar cada bloque en la diapositiva mientras lo enuncia. Esto orienta al comite
sobre la estructura y establece expectativas claras.
\end{tipbox}

\vspace{8pt}

% ==============================================================
%  DIAPOSITIVA 3
% ==============================================================
\subsection{S03 --- Antecedentes}

\begin{timebox}
Tiempo estimado: 0:45 min \quad|\quad Acumulado: 1:35 min
\end{timebox}
\vspace{4pt}

\begin{slidebox}{Guion}
``?`Por que necesitamos una alternativa a los metodos clasicos? Simular speckle
optico requiere resolver la ecuacion de Helmholtz para miles de perfiles de rugosidad
distintos. \textbf{FEM} exige una malla con resolucion menor que la longitud de onda:
para $\lambda = 638$ nm en un dominio 2D ya hablamos de millones de incognitas por
realizacion, y cada nueva superficie requiere construir la malla desde cero. \textbf{Monte
Carlo} necesita rastrear un numero masivo de fotones individuales, con un costo muy
alto por realizacion.
\vspace{4pt}

Ambos metodos son correctos pero intratables para estudios estadisticos que requieren
cientos o miles de realizaciones. Las \textbf{PINNs-SIREN} ofrecen una alternativa
libre de malla, con inferencia en milisegundos una vez entrenada la red.''
\end{slidebox}
\vspace{4pt}

\begin{tipbox}
La diapositiva tiene \textbf{dos columnas}: primero senalar la izquierda (problemas de FEM
y Monte Carlo), luego la derecha (ventajas de PINN). El contraste visual refuerza el
mensaje.
\end{tipbox}

\vspace{8pt}

% ==============================================================
%  DIAPOSITIVA 4
% ==============================================================
\subsection{S04 --- Planteamiento del Problema}

\begin{timebox}
Tiempo estimado: 0:45 min \quad|\quad Acumulado: 2:20 min
\end{timebox}
\vspace{4pt}

\begin{slidebox}{Guion}
``El problema central es la escalabilidad. Para estudios estadisticos de speckle
necesitamos resolver el campo electromagnetico para cientos de superficies rugosas
distintas. FEM tiene tres limitaciones criticas: requiere alta resolucion de malla
(mayor uso de RAM), tiempo de ejecucion elevado por realizacion, y debe resolver
desde cero ante cualquier cambio en la condicion de frontera.

La pregunta que surge de esto es: ?`podemos entrenar una red neuronal que
\textbf{aprenda la fisica de Helmholtz una vez} y luego evalue nuevas condiciones
de frontera en milisegundos sin repetir el entrenamiento completo?''
\end{slidebox}
\vspace{4pt}

\begin{tipbox}
\textbf{Enfatizar} ``una vez'' y ``milisegundos'' como la propuesta de valor central.
Pausar brevemente despues para que el mensaje se asiente.
\end{tipbox}

\vspace{8pt}

% ==============================================================
%  DIAPOSITIVA 5
% ==============================================================
\subsection{S05 --- Preguntas de Investigacion}

\begin{timebox}
Tiempo estimado: 0:30 min \quad|\quad Acumulado: 2:50 min
\end{timebox}
\vspace{4pt}

\begin{slidebox}{Guion}
``De ahi se derivan dos preguntas de investigacion concretas y verificables.
\textbf{Primera:} ?`puede una PINN-SIREN aproximar la solucion de la ecuacion de
Helmholtz 2D con error $L^2$ menor al 5\%?
\textbf{Segunda:} ?`puede esa misma arquitectura simular un patron de speckle con
contraste $C \approx 1$, que es el criterio de Goodman para speckle completamente
desarrollado?
Al final de la presentacion veremos la respuesta a la primera pregunta; la segunda
esta en progreso con NB03.''
\end{slidebox}
\vspace{4pt}

\begin{tipbox}
Senalar cada pregunta en la diapositiva. Hacer \textbf{pausa} para que el comite lea.
El numero ``5\%'' y ``$C \approx 1$'' deben quedar en la memoria del comite.
\end{tipbox}

\vspace{8pt}

% ==============================================================
%  DIAPOSITIVA 6
% ==============================================================
\subsection{S06 --- Hipotesis}

\begin{timebox}
Tiempo estimado: 0:30 min \quad|\quad Acumulado: 3:20 min
\end{timebox}
\vspace{4pt}

\begin{slidebox}{Guion}
``La hipotesis central es que si es posible. Una PINN con activacion sinusoidal
puede resolver la ecuacion de Helmholtz 2D con campo complejo, alcanzando un error
$L^2$ menor al 5\% y un contraste de speckle $C$ dentro del 10\% del valor teorico de
$1$. Los resultados que presentare a continuacion \textbf{confirman la hipotesis para
los casos 1D y 2D}; NB03 completara la validacion para el caso de speckle.''
\end{slidebox}
\vspace{4pt}

\begin{tipbox}
Ser conciso. Esta diapositiva es un puente entre el problema y los resultados.
No extenderse; el comite ya entendio el contexto.
\end{tipbox}

\vspace{8pt}

% ==============================================================
%  DIAPOSITIVA 7
% ==============================================================
\subsection{S07 --- Objetivos de Investigacion}

\begin{timebox}
Tiempo estimado: 0:30 min \quad|\quad Acumulado: 3:50 min
\end{timebox}
\vspace{4pt}

\begin{slidebox}{Guion}
``El objetivo general es \textbf{desarrollar y validar una PINN-SIREN para
la simulacion de speckle optico}. Los objetivos especificos son: primero,
validar la arquitectura en 1D con solucion analitica conocida; segundo, extender
a 2D con campo complejo; tercero, implementar metricas estadisticas de speckle
---contraste, distribucion de intensidad---; y cuarto, comparar el tiempo de
calculo contra el estado del arte.''
\end{slidebox}
\vspace{4pt}

\begin{tipbox}
Leer los objetivos especificos directamente de la slide es aceptable aqui.
Son referencia para que el comite evalue el avance. Ir despacio.
\end{tipbox}

\vspace{8pt}

% ==============================================================
%  DIAPOSITIVA 8
% ==============================================================
\subsection{S08 --- Cronograma}

\begin{timebox}
Tiempo estimado: 0:25 min \quad|\quad Acumulado: 4:15 min
\end{timebox}
\vspace{4pt}

\begin{slidebox}{Guion}
``En cuanto al cronograma: los avances hasta mayo estan completados. NB01
---Helmholtz 1D--- y NB02 ---campo complejo 2D--- ya fueron validados y son
los resultados que presentare en esta sesion. El trabajo pendiente para el
segundo semestre es NB03 ---simulacion de speckle propiamente dicho--- y NB04
---benchmark formal contra FEniCSx.''
\end{slidebox}
\vspace{4pt}

\begin{tipbox}
Senalar en la tabla los meses completados (color verde o marcados) y los pendientes.
Esto demuestra progreso sistematico y planificacion.
\end{tipbox}

\vspace{8pt}

% ==============================================================
%  DIAPOSITIVA 9
% ==============================================================
\subsection{S09 --- Motivacion}

\begin{timebox}
Tiempo estimado: 0:30 min \quad|\quad Acumulado: 4:45 min
\end{timebox}
\vspace{4pt}

\begin{slidebox}{Guion}
``?`Por que simular speckle es relevante? El patron de speckle aparece en
sistemas opticos con luz coherente: metrologia optica, deteccion de defectos
superficiales, interferometria, diagnostico medico con laser ---como en tecnicas
de imagen retiniana---, LIDAR y comunicaciones de fibra optica. En todos estos
dominios se necesita predecir la estadistica del patron de speckle para disenar
y calibrar los sistemas. La simulacion computacional acelerada es por tanto un
habilitador tecnologico de alto impacto.''
\end{slidebox}
\vspace{4pt}

\begin{tipbox}
Si el tiempo lo permite, mencionar una aplicacion concreta relacionada con el
contexto del comite. Esta slide es de ``venta'' del problema: mostrar entusiasmo.
\end{tipbox}

\vspace{8pt}

% ==============================================================
%  DIAPOSITIVA 10
% ==============================================================
\subsection{S10 --- Marco Teorico: Helmholtz y Speckle}

\begin{timebox}
Tiempo estimado: 1:00 min \quad|\quad Acumulado: 5:45 min
\end{timebox}
\vspace{4pt}

\begin{slidebox}{Guion}
``El fundamento fisico del trabajo es la \textbf{ecuacion de Helmholtz escalar 2D}:
\[
  \nabla^2 E(\mathbf{r}) + k^2 E(\mathbf{r}) = 0
\]
donde $E$ es el campo electrico complejo, $k = 2\pi/\lambda$ es el numero de onda.
En nuestro caso usamos un dominio adimensional $[0,1]^2$ con $k = 2\pi$, equivalente
a un laser diodo rojo de 638 nm. En coordenadas cartesianas el laplaciano es:
\[
  \nabla^2 E = \frac{\partial^2 E}{\partial x^2} + \frac{\partial^2 E}{\partial y^2}
\]
El speckle surge cuando la superficie reflectante tiene rugosidad mayor que $\lambda$,
de modo que cada punto refleja la onda con una fase aleatoria. La interferencia
constructiva y destructiva de esas contribuciones genera el patron de speckle.''
\end{slidebox}
\vspace{4pt}

\begin{tipbox}
\textbf{Senalar la ecuacion} en la diapositiva mientras la enuncia. El comite es
de CS, no de fisica: explicar brevemente que $k$ es el numero de onda y que
el campo es complejo. Pronunciar ``nabla cuadrado'' o ``laplaciano de E''.
\end{tipbox}

\vspace{8pt}

% ==============================================================
%  DIAPOSITIVA 11
% ==============================================================
\subsection{S11 --- PINNs y Arquitectura SIREN}

\begin{timebox}
Tiempo estimado: 1:00 min \quad|\quad Acumulado: 6:45 min
\end{timebox}
\vspace{4pt}

\begin{slidebox}{Guion}
``Las \textbf{Physics-Informed Neural Networks} incorporan la ecuacion diferencial
directamente en la funcion de perdida. En lugar de datos etiquetados, el residuo
de Helmholtz actua como supervision:
\[
  \mathcal{L} = \mathcal{L}_\text{BC} + \lambda_\text{fis} \cdot \mathcal{L}_\text{fisica}
\]
La clave de nuestra propuesta es la activacion \textbf{SIREN}: $\phi(z) = \sin(\omega_0 \cdot z)$.
?`Por que seno en lugar de ReLU o tanh? Porque Helmholtz requiere derivadas de segundo
orden. ReLU tiene segunda derivada cero ---completamente inutil para el laplaciano---.
tanh se satura: sus derivadas de orden alto se aproximan a cero. SIREN tiene derivadas
de cualquier orden, todas son sinusoides, lo que hace la perdida de fisica diferenciable
y bien condicionada.''
\end{slidebox}
\vspace{4pt}

\begin{tipbox}
El \textbf{punto critico para el comite de CS}: las derivadas $\partial^2 E/\partial x^2$
se calculan con \code{torch.autograd.grad()} encadenado dos veces --- diferenciacion
automatica exacta, sin discretizacion. Mencionarlo brevemente aqui y profundizar en S14.
\end{tipbox}

\vspace{8pt}

% ==============================================================
%  DIAPOSITIVA 12
% ==============================================================
\subsection{S12 --- Arquitectura SIREN 5x128}

\begin{timebox}
Tiempo estimado: 1:00 min \quad|\quad Acumulado: 7:45 min
\end{timebox}
\vspace{4pt}

\begin{slidebox}{Guion}
``La arquitectura concreta para el caso 2D es: \textbf{5 capas de 128 neuronas}
con activacion seno. Entrada: $(x, y)$ ---dos neuronas---. Salida: $(E_\text{real},
E_\text{imag})$ ---dos neuronas---. Total: 66,690 parametros.

La \textbf{inicializacion de pesos} es especial y critica. La capa de entrada se
inicializa con $W_0 \sim U(-1/n, +1/n)$. Las capas ocultas con $W_i \sim U(-\sqrt{6/n}/\omega_0,
+\sqrt{6/n}/\omega_0)$. Esta inicializacion garantiza que las activaciones sinusoidales
no se saturen desde el inicio y que el espectro de frecuencias sea estable a traves de
todas las capas.

La \textbf{estrategia de entrenamiento} es bifasica: primero Adam para exploracion
global, luego L-BFGS con busqueda de linea strong Wolfe para refinamiento fino.''
\end{slidebox}
\vspace{4pt}

\begin{tipbox}
El diagrama de la arquitectura en la slide muestra la conexion capa a capa.
Senalarlo brevemente. Si alguien pregunta por el numero de parametros:
$5 \times (128 \times 128) + $ capas entrada/salida $\approx 66{,}690$.
\end{tipbox}

\vspace{8pt}

% ==============================================================
%  DIAPOSITIVA 13
% ==============================================================
\subsection{S13 --- Metodologia: Muestreo, Perdida, Optimizacion}

\begin{timebox}
Tiempo estimado: 0:45 min \quad|\quad Acumulado: 8:30 min
\end{timebox}
\vspace{4pt}

\begin{slidebox}{Guion}
``Los tres pilares metodologicos del experimento. \textbf{Primero: muestreo LHS.}
Latin Hypercube Sampling garantiza que los 3,000 puntos de colocacion cubran el
dominio $[0,1]^2$ de manera uniforme: ni clustering ni zonas vacias. Esto es
superior al muestreo aleatorio simple en 2D.

\textbf{Segundo: funcion de perdida PINN} con $\lambda = 0.1$, calibrado
experimentalmente para evitar desbordamiento de gradientes en float32.

\textbf{Tercero: optimizacion bifasica.} Adam explora el paisaje de perdida
durante 15,000 epocas con \emph{early stopping}. Una vez que Adam converge,
L-BFGS usa la curvatura del Hessiano aproximado para refinamiento de alta
precision.''
\end{slidebox}
\vspace{4pt}

\begin{tipbox}
La slide tiene tres columnas, una por pilar. Senalar cada una mientras la enuncia.
El comite puede preguntar por el valor $\lambda = 0.1$; la respuesta detallada
esta en la diapositiva de respaldo \textbf{S32}.
\end{tipbox}

\vspace{8pt}

% ==============================================================
%  DIAPOSITIVA 14
% ==============================================================
\subsection{S14 --- Pipeline de Implementacion}

\begin{timebox}
Tiempo estimado: 0:45 min \quad|\quad Acumulado: 9:15 min
\end{timebox}
\vspace{4pt}

\begin{slidebox}{Guion}
``Aqui ven el pipeline real del codigo implementado en seis etapas.
\textbf{Etapa 1:} muestreo LHS en \code{xy} $\in [0,1]^2$.
\textbf{Etapa 2:} forward pass a traves de \code{PINN\_2D\_SIREN.forward(xy)} en
\code{models.py}.
\textbf{Etapa 3:} calculo del residuo de Helmholtz en \code{losses.py} ---
\code{autograd.grad(E, xy)} nos da el gradiente de $E$, y encadenamos una segunda
llamada para obtener el laplaciano exacto: $\nabla^2 E$. Sin diferencias finitas,
sin malla. El resultado es exacto a la precision de float32.
\textbf{Etapa 4:} computo de la funcion de perdida.
\textbf{Etapa 5:} optimizacion bifasica en \code{training.py}.
\textbf{Etapa 6:} evaluacion de metricas $L^2$, $R^2$ y persistencia del modelo
con \code{utils.py}.''
\end{slidebox}
\vspace{4pt}

\begin{tipbox}
Esta es la diapositiva que mas valora el \textbf{comite computacional}.
Enfatizar la Etapa 3: la diferenciacion automatica encadenada es la innovacion
de implementacion. Mostrar mentalmente el codigo: \code{autograd.grad(grad\_f, xy)}.
\end{tipbox}

\vspace{8pt}

% ==============================================================
%  DIAPOSITIVA 15
% ==============================================================
\subsection{S15 --- Arquitectura del Codigo}

\begin{timebox}
Tiempo estimado: 0:30 min \quad|\quad Acumulado: 9:45 min
\end{timebox}
\vspace{4pt}

\begin{slidebox}{Guion}
``El codigo esta organizado en cuatro modulos especializados. \code{models.py}
define la arquitectura SIREN con la inicializacion especial. \code{losses.py}
implementa el residuo de Helmholtz con autograd y la funcion de perdida compuesta.
\code{training.py} contiene el loop de entrenamiento Adam mas L-BFGS con el
patron de \emph{closure}. \code{utils.py} maneja las metricas y la persistencia
de modelos en formato \code{.pt}. Los notebooks NB01 y NB02 consumen estos
modulos y exportan modelos reutilizables.''
\end{slidebox}
\vspace{4pt}

\begin{tipbox}
Ir rapido aqui; el pipeline anterior ya dio el contexto. Esta slide es un
mapa de navegacion del codigo. Si el comite pregunta sobre el diseno:
``la modularizacion permite reutilizar pesos entre notebooks sin reentrenar.''
\end{tipbox}

\vspace{8pt}

% ==============================================================
%  DIAPOSITIVAS 16-17 (par de animacion)
% ==============================================================
\subsection{S16/S17 --- NB01: Validacion Helmholtz 1D}

\begin{timebox}
Tiempo estimado: 1:00 min \quad|\quad Acumulado: 10:45 min
\end{timebox}
\vspace{4pt}

\begin{warningbox}
S16 y S17 son el mismo contenido en dos estados de animacion. \textbf{Presentar
solo desde S17} (la version completa con todos los elementos visibles).
Si usa la version animada, avanzar S16 rapidamente y detenerse en S17.
\end{warningbox}
\vspace{4pt}

\begin{slidebox}{Guion}
``Para validar el enfoque, comenzamos con el caso mas simple: Helmholtz 1D con
solucion analitica conocida, $E_\text{exacta}(x) = \cos(kx)$, con condiciones
de frontera $E(0) = E(1) = 1$ y tres puntos intermedios conocidos.

La red SIREN 5$\times$64 ---16,833 parametros--- converge en 251 segundos.
Los resultados son: \textbf{error $L^2 = 0.006\%$} y \textbf{$R^2 = 1.000000$}
con seis cifras significativas. En la grafica, la curva azul de la PINN es
practicamente indistinguible de la solucion analitica roja. Esta precision
supera ampliamente la meta del 5\%.''
\end{slidebox}
\vspace{4pt}

\begin{tipbox}
\textbf{Estos numeros son su logro mas fuerte}. Decirlos con seguridad y pausar
para que el comite los asimile. Senalar la grafica de comparacion en la slide.
``0.006\% de error L2, R cuadrado de 1.0, seis cifras decimales.''
\end{tipbox}

\vspace{8pt}

% ==============================================================
%  DIAPOSITIVAS 18-19 (par de animacion)
% ==============================================================
\subsection{S18/S19 --- NB01: Analisis de Error y Convergencia}

\begin{timebox}
Tiempo estimado: 0:45 min \quad|\quad Acumulado: 11:30 min
\end{timebox}
\vspace{4pt}

\begin{warningbox}
S18 y S19 son par de animacion. \textbf{Presentar desde S19} (completa).
\end{warningbox}
\vspace{4pt}

\begin{slidebox}{Guion}
``El analisis de error confirma la precision. El histograma de error absoluto
esta concentrado cerca de cero con distribucion simetrica ---sin sesgos
sistematicos---. La correlacion puntual con la solucion analitica muestra
$R^2 = 1$. La curva de convergencia tiene la forma caracteristica de la
estrategia bifasica: descenso rapido durante Adam y refinamiento abrupto
al iniciar L-BFGS, que convergio en 202 iteraciones.''
\end{slidebox}
\vspace{4pt}

\begin{tipbox}
Ir rapido; el comite ya entendio el resultado clave en S17. Esta slide
refuerza el mensaje con evidencia grafica. Mencionar la ``firma'' de L-BFGS
en la curva de convergencia ---el salto brusco al final--- pues demuestra que
la estrategia bifasica funciono.
\end{tipbox}

\vspace{8pt}

% ==============================================================
%  DIAPOSITIVAS 20-23 (cuadruplete de animacion)
% ==============================================================
\subsection{S20--S23 --- NB02: Helmholtz 2D con Campo Complejo}

\begin{timebox}
Tiempo estimado: 1:30 min \quad|\quad Acumulado: 13:00 min
\end{timebox}
\vspace{4pt}

\begin{warningbox}
S20--S23 son cuatro estados de animacion del mismo contenido.
\textbf{Detenerse en S23} (la version con todas las figuras visibles).
Si usa version animada, avanzar S20--S22 rapidamente.
\end{warningbox}
\vspace{4pt}

\begin{slidebox}{Guion}
``El caso 2D es significativamente mas desafiante: debemos predecir simultaneamente
la parte real e imaginaria del campo complejo. La solucion exacta es una onda plana:
$E_\text{exacta} = e^{i(k_x x + k_y y)}$.

La red SIREN 5$\times$128 con 66,690 parametros alcanza:
\begin{itemize}
  \item \textbf{Error $L^2$ promedio: 0.171\%} (0.214\% real, 0.127\% imaginaria)
  \item \textbf{Tiempo de entrenamiento: 299 segundos}
  \item \textbf{L-BFGS: 1,035 iteraciones}
\end{itemize}
En la figura pueden ver el campo de amplitud $|E|$ y la fase $\angle E$ predichos
por la PINN: visualmente indistinguibles de la solucion analitica. Los errores
puntuales son del orden de $10^{-3}$, distribuidos uniformemente en el dominio.''
\end{slidebox}
\vspace{4pt}

\begin{tipbox}
\textbf{Enfatizar que el campo es complejo} ---parte real e imaginaria---.
Muchos metodos de ML solo trabajan con reales. SIREN maneja ambas salidas
simultaneamente con una sola red. Senalar el mapa de error en la figura:
el error es homogeneo, sin zonas de alta concentracion.
\end{tipbox}

\vspace{8pt}

% ==============================================================
%  DIAPOSITIVA 24
% ==============================================================
\subsection{S24 --- NB02: Analisis Multi-Semilla}

\begin{timebox}
Tiempo estimado: 0:40 min \quad|\quad Acumulado: 13:40 min
\end{timebox}
\vspace{4pt}

\begin{slidebox}{Guion}
``Para validar la robustez estadistica del metodo, entrenamos con tres semillas
distintas: 42, 123 y 777. Los resultados son:
\[
  L^2_\text{promedio} = 0.155 \pm 0.020\%
\]
La varianza entre semillas es minima: la desviacion estandar del 0.020\% es dos
ordenes de magnitud menor que el umbral del 5\%. Esto indica que el resultado no
es un accidente de inicializacion sino que la arquitectura SIREN converge de forma
consistente independientemente de los pesos iniciales.''
\end{slidebox}
\vspace{4pt}

\begin{tipbox}
Este resultado es clave para la solidez cientifica. El comite puede preguntar:
``?`como saben que no es sobreajuste?'' Respuesta: la solucion exacta es conocida
y los 3 seeds dan el mismo resultado.
\end{tipbox}

\vspace{8pt}

% ==============================================================
%  DIAPOSITIVA 25
% ==============================================================
\subsection{S25 --- Comparativa con el Estado del Arte}

\begin{timebox}
Tiempo estimado: 0:45 min \quad|\quad Acumulado: 14:25 min
\end{timebox}
\vspace{4pt}

\begin{slidebox}{Guion}
``La comparativa con el estado del arte es la contribucion clave del avance.
Schoder y Kraxberger (2024) reportaron PINNs convencionales con activacion tanh
para propagacion de ondas. Nuestros resultados muestran una aceleracion de:
\begin{itemize}
  \item \textbf{$\times$415} en el caso 1D
  \item \textbf{$\times$11.2} en el caso 2D
\end{itemize}
Es importante ser precisos: la comparacion es \textbf{indicativa, no directa}.
Schoder trabajaron en 3D y nosotros en 2D. Sin embargo, el resultado establece
claramente el potencial de SIREN sobre activaciones convencionales para este tipo
de problemas de onda.''
\end{slidebox}
\vspace{4pt}

\begin{tipbox}
\textbf{Mencionar proactivamente la salvedad de 3D vs 2D.} Esto demuestra honestidad
cientifica y previene que el comite lo use como critica. Decir algo como: ``somos
conscientes de que la comparacion es en distinto numero de dimensiones; por eso
NB04 hara el benchmark directo contra FEniCSx en las mismas condiciones.''
\end{tipbox}

\vspace{8pt}

% ==============================================================
%  DIAPOSITIVA 26
% ==============================================================
\subsection{S26 --- Discusion: Hallazgos Principales}

\begin{timebox}
Tiempo estimado: 0:40 min \quad|\quad Acumulado: 15:05 min
\end{timebox}
\vspace{4pt}

\begin{slidebox}{Guion}
``Resumiendo los hallazgos de este avance: NB01 alcanzo $L^2 = 0.006\%$ con
$R^2 = 1.000000$, superando ampliamente la meta del 5\%. NB02 alcanzo $L^2 = 0.171\%$
para campo complejo 2D, igualmente muy por debajo del umbral. La validacion
multi-semilla confirma robustez estadistica con $0.155 \pm 0.020\%$.

La hipotesis de que PINN-SIREN puede resolver Helmholtz con la precision requerida
para speckle \textbf{queda respaldada por los datos}. El siguiente paso es NB03,
donde se sustituye la condicion de frontera de onda plana por fase aleatoria para
generar el patron de speckle.''
\end{slidebox}
\vspace{4pt}

\begin{tipbox}
Esta es la slide de cierre de resultados. Hablar con confianza. Los numeros
hablan por si solos; no agregar demasiada narrativa.
\end{tipbox}

\vspace{8pt}

% ==============================================================
%  DIAPOSITIVA 27
% ==============================================================
\subsection{S27 --- Trabajo Futuro y Proximos Pasos}

\begin{timebox}
Tiempo estimado: 0:30 min \quad|\quad Acumulado: 15:35 min
\end{timebox}
\vspace{4pt}

\begin{slidebox}{Guion}
``Los proximos pasos en orden de prioridad son: \textbf{NB03}, simulacion de
speckle con condicion de frontera de fase aleatoria $\phi \sim U(0, 2\pi)$ en
$y=0$; el criterio de exito es contraste $C \approx 1$ y distribucion de intensidad
exponencial. \textbf{NB04}, benchmark formal contra FEniCSx bajo las mismas
condiciones para medir el factor de aceleracion $S = T_\text{FEM}/T_\text{PINN}$.
Estos dos experimentos completan la hipotesis central de la tesis.''
\end{slidebox}
\vspace{4pt}

\begin{tipbox}
Si alguien pregunta por el calendario: NB03 se estima para julio-agosto 2026,
NB04 para agosto-septiembre. La defensa se proyecta para finales de 2026.
\end{tipbox}

\vspace{8pt}

% ==============================================================
%  DIAPOSITIVA 28
% ==============================================================
\subsection{S28 --- Resumen del Avance}

\begin{timebox}
Tiempo estimado: 0:30 min \quad|\quad Acumulado: 16:05 min
\end{timebox}
\vspace{4pt}

\begin{slidebox}{Guion}
``En sintesis, en este segundo avance: se demostro precision $L^2 < 0.5\%$ en
ambos casos experimentados ---muy por debajo de la meta del 5\%---. La arquitectura
SIREN es estable y reproducible con multiples semillas. El pipeline modular esta
implementado y documentado. Y los fundamentos tecnicos y computacionales estan
sentados para abordar NB03 y NB04 en el siguiente semestre.''
\end{slidebox}
\vspace{4pt}

\begin{tipbox}
Enunciar la lista de logros con claridad y en voz firme. Esta es la ``venta final''
del avance antes del turno de preguntas.
\end{tipbox}

\vspace{8pt}

% ==============================================================
%  DIAPOSITIVA 29
% ==============================================================
\subsection{S29 --- Referencias}

\begin{timebox}
Tiempo estimado: 0:10 min \quad|\quad Acumulado: 16:15 min
\end{timebox}
\vspace{4pt}

\begin{slidebox}{Guion}
``Las referencias clave son: Raissi et al. 2019 ---el trabajo fundacional de PINNs
en \emph{Journal of Computational Physics}---, Sitzmann et al. 2020 ---SIREN en
NeurIPS---, y Schoder y Kraxberger 2024 para la comparativa. La bibliografia
completa esta en la diapositiva y en la tesis.''
\end{slidebox}
\vspace{4pt}

\begin{tipbox}
\textbf{No leer} todas las referencias. Mencionar solo las tres clave.
El comite puede verificar la diapositiva por su cuenta.
\end{tipbox}

\vspace{8pt}

% ==============================================================
%  DIAPOSITIVA 30
% ==============================================================
\subsection{S30 --- Cierre y Agradecimiento}

\begin{timebox}
Tiempo estimado: 0:20 min \quad|\quad Acumulado: \textbf{16:35 min} \quad|\quad Margen: $\approx$3:25 min
\end{timebox}
\vspace{4pt}

\begin{slidebox}{Guion}
``Muchas gracias por su atencion y por el tiempo dedicado a evaluar este avance.
Quedo a sus ordenes para responder las preguntas del comite.''
\end{slidebox}
\vspace{4pt}

\begin{tipbox}
Sonrisa, contacto visual con cada sinodal. Postura abierta. No agregar
mas contenido ni justificaciones no solicitadas. Esperar las preguntas
con calma.

\textbf{Recuerde:} las diapositivas de respaldo S31--S33 estan disponibles
para preguntas sobre ablacion $\lambda$ e inicializacion SIREN. No es necesario
mencionarlas; mostrarlas solo si se pregunta.
\end{tipbox}

\vspace{8pt}

% ==============================================================
%  DIAPOSITIVAS DE RESPALDO (no presentar salvo pregunta)
% ==============================================================
\subsection{S31--S33 --- Diapositivas de Respaldo (no presentar)}

\begin{warningbox}
\textbf{Estas slides NO forman parte del guion principal.}
Mostrarlas solo si el comite pregunta sobre:
\begin{itemize}
  \item Ablacion de $\lambda$ (por que 0.1 y no 0.01 o 1.0) $\to$ \textbf{S32}
  \item Inicializacion de pesos SIREN (Sitzmann et al.) $\to$ \textbf{S33}
  \item Diapositiva divisor de respaldo $\to$ \textbf{S31}
\end{itemize}
El contenido de estas slides esta cubierto en detalle en la Seccion 3
(Preguntas y Respuestas) de este documento.
\end{warningbox}

\vspace{16pt}

% --- Tabla resumen final ---
\subsection{Tabla Resumen de Tiempos}

\begin{table}[h!]
\centering
\small
\begin{tabular}{@{}clccc@{}}
\toprule
\textbf{Slide(s)} & \textbf{Titulo} & \textbf{Tiempo} & \textbf{Acumulado} & \textbf{Bloque} \\
\midrule
S01 & Portada & 0:25 & 0:25 & I \\
S02 & Contenido & 0:25 & 0:50 & I \\
S03 & Antecedentes & 0:45 & 1:35 & I \\
S04 & Planteamiento & 0:45 & 2:20 & I \\
S05 & Preguntas & 0:30 & 2:50 & I \\
S06 & Hipotesis & 0:30 & 3:20 & I \\
S07 & Objetivos & 0:30 & 3:50 & I \\
S08 & Cronograma & 0:25 & 4:15 & I \\
S09 & Motivacion & 0:30 & 4:45 & I \\
\midrule
S10 & Marco Teorico Helmholtz & 1:00 & 5:45 & II \\
S11 & PINNs y SIREN & 1:00 & 6:45 & II \\
S12 & Arquitectura SIREN 5$\times$128 & 1:00 & 7:45 & II \\
S13 & Metodologia & 0:45 & 8:30 & II \\
S14 & Pipeline & 0:45 & 9:15 & II \\
S15 & Arquitectura codigo & 0:30 & 9:45 & II \\
\midrule
S16/17 & NB01 Validacion 1D & 1:00 & 10:45 & III \\
S18/19 & NB01 Analisis error & 0:45 & 11:30 & III \\
S20--23 & NB02 Helmholtz 2D & 1:30 & 13:00 & III \\
S24 & NB02 Multi-semilla & 0:40 & 13:40 & III \\
\midrule
S25 & Comparativa estado del arte & 0:45 & 14:25 & IV \\
S26 & Discusion hallazgos & 0:40 & 15:05 & IV \\
S27 & Trabajo futuro & 0:30 & 15:35 & IV \\
S28 & Resumen & 0:30 & 16:05 & IV \\
S29 & Referencias & 0:10 & 16:15 & IV \\
S30 & Cierre & 0:20 & 16:35 & IV \\
\midrule
\textbf{Total} & & \textbf{16:35} & & \\
\textbf{Margen} & & \textbf{3:25} & & \\
\bottomrule
\end{tabular}
\caption{Resumen de tiempos por diapositiva. El margen de $\approx$3.5 min absorbe
preguntas aclaratorias durante la presentacion o transiciones lentas.}
\end{table}
